\chapter{変数定義と採用根拠：四要素モデルの実証的写像}
\label{chap:variables}

第4章では，Merton モデルおよび Bharath \& Shumway（2008）（以下 B\&S）の
multiplicative → log → linear の構造を基礎として，
新電力企業の倒産構造を反映した潜在倒産距離 $Z_{\text{lin}}$ を理論的に構築した。
本章では，この理論モデルを実証分析に接続するために，
$Z_{\text{lin}}$ の構成要素である収益性 $\mu$，
外部ショック感応度 $\sigma$，支援能力 $S$，倒産閾値 $D$ を，
観測可能な変数として定義し，その採用根拠を整理する。

新電力企業の倒産は，外部ショックが短期間で資金繰りに波及し，
複数の要因が掛け算的に悪化することで急激に顕在化するという特徴を持つ。
この構造は Merton 型モデルの multiplicative な性質と整合的であり，
また B\&S が示した「理論構造を観測可能な proxy によって離散時間に写像する」
という方法論は，新電力企業の倒産リスクを実証的に扱ううえで有効である。

本章では，まず四要素モデルの各構成要素が
Merton/B\&S の理論構造のどの部分に対応するかを明確化し（5.1 節），
続いて収益性 $\mu$（5.2 節），外部ショック感応度 $\sigma$（5.3 節），
支援能力 $S$（5.4 節），倒産閾値 $D$（5.5 節）について，
変数定義と採用根拠を順に整理する。
これにより，第6章で構築するロジットモデルにおいて，
$Z_{\text{lin}}$ を実証的に推定可能な形で扱うための基盤を整える。


% ---------------------------------------------
% 5.1 本章の目的と基本方針
% 本章では、理論モデル Z_lin の構成要素を
% 実証で測定可能な変数として定義する。
% 理論 → データ への写像を明確化する章である。
% ---------------------------------------------

\section{本章の目的と基本方針}

第4章では，Merton モデルおよび Bharath \& Shumway（2008）（以下 B\&S）の
multiplicative → log → linear の構造を基礎として，
新電力企業の倒産構造を反映した潜在倒産距離 $Z_{\text{lin}}$ を理論的に構築した。
本章では，この $Z_{\text{lin}}$ を実証分析に接続するために，
その構成要素である収益性 $\mu$，外部ショック感応度 $\sigma$，
支援能力 $S$，倒産閾値 $D$ を，観測可能な変数として定義する。

理論モデルを実証に写像する際には，B\&S が行ったように，
理論的パラメータを観測可能な proxy によって置き換える必要がある。
すなわち，Merton モデルにおける $\mu$ や $\sigma_V$ が
株価データによって proxy 化されたのと同様に，
本研究でも新電力企業の制度的特徴を踏まえた proxy 変数を構築する。

本章の目的は，以下の 3 点である。

\begin{enumerate}
    \item $Z_{\text{lin}}$ の構成要素（$\mu, \sigma, S, D$）を実証的に測定可能な変数として定義すること
    \item 各変数の採用根拠を，Merton/B\&S の理論構造および新電力市場の制度的特徴に基づき明確化すること
    \item 第6章で構築するロジットモデルにおいて，これらの変数がどのように倒産確率に影響するかの理論的期待値（符号条件）を整理すること
\end{enumerate}

以上により，本章は「理論モデル（第4章）→ 実証モデル（第6章）」をつなぐ
橋渡しとして位置づけられる。
次節以降では，四要素モデルの各構成要素について，
変数定義と採用根拠を順に整理する。

% ---------------------------------------------
% 5.2 収益性 μ の定義と採用根拠（修正版）
% ---------------------------------------------

\section{収益性 $\mu$ の定義と採用根拠}

本節では，潜在倒産距離 $Z_{\text{lin}}$ の構成要素のうち，
収益性 $\mu$ を実証的にどのように測定するかを定義し，
その採用根拠を Merton/B\&S の理論構造および新電力市場の制度的特徴に基づき整理する。

\subsection{理論的対応：Merton/B\&S における $\mu$ の位置づけ}

Merton モデルにおける $\mu$ は企業価値 $V_t$ の期待成長率であり，
企業の基礎的な収益力を表す潜在パラメータである。
しかし企業価値 $V_t$ は観測不可能であるため，
Bharath \& Shumway（2008）（以下 B\&S）は，
$\mu$ を株価リターンの平均など観測可能な proxy によって置き換えた。

すなわち，



\[
\mu_{\text{theory}} \approx \hat{\mu}_{\text{proxy}}
\]



という形で，理論上の収益性を実証的に扱える形に変換している。
本研究もこの proxy 化の思想を継承し，
新電力企業の制度的特徴に適合する形で $\mu$ を定義する。

\subsection{新電力市場における収益性の意味}

新電力企業の収益性は，外部ショックに対する耐性を決定する基本要因である。
具体的には，

\begin{itemize}
    \item 収益性が高い企業ほど，JEPX 価格急騰や制度コスト増加に対して資金繰り余力が大きい
    \item 収益性が低い企業は，外部ショック発生時に即座に赤字転落しやすい
    \item 収益性は倒産距離を押し上げる（$a>0$）
\end{itemize}

という特徴を持つ。
したがって，収益性 $\mu$ は新電力企業の倒産構造において
「基礎体力」を表す中心的な要因である。

\subsection{本来の定義：営業利益率（理論上の $\mu$）}

理論的には，収益性 $\mu$ は企業の本業収益力を表す指標として，
営業利益率を用いることが最も自然である。



\[
\mu_{\text{theory}} = \frac{\text{営業利益}}{\text{売上高}}
\]



営業利益率は，

\begin{itemize}
    \item 調達費用の変動が直接反映される
    \item 企業間比較が容易である
    \item 倒産研究で標準的に用いられる
\end{itemize}

という理由から，理論上の $\mu$ の定義として妥当である。

\subsection{データ制約：倒産企業は損益計算書を公表しない}

しかし，日本の決算公告制度では \textbf{貸借対照表のみの公表が認められており}，
倒産企業の多くは損益計算書（売上高・営業利益）を公表しない。
したがって，営業利益率を用いた $\mu_{\text{theory}}$ を
実証的に測定することは困難である。

この制度的制約は，本研究に固有の問題ではなく，
倒産研究一般に共通する構造的制約である。

\subsection{実証上の定義：貸借対照表ベースの proxy（$\mu$ の代替指標）}

この制約を踏まえ，本研究では B\&S の proxy 化の思想に従い，
収益性 $\mu$ を貸借対照表ベースの指標によって proxy 化する。
具体的には，以下の 3 つの指標を候補とする。

\noindent\textbf{（1）純資産の増減率（最も理論的整合性が高い）}



\[
\mu_1 = \frac{\Delta \text{純資産}}{\text{総資産}}
\]



純資産の増減は損益の蓄積を反映するため，
営業利益率の自然な代替指標となる。

\noindent\textbf{（2）純資産比率（財務健全性の proxy）}



\[
\mu_2 = \frac{\text{純資産}}{\text{総資産}}
\]



資本の厚みは外部ショックに対する耐性を表す。

\noindent\textbf{（3）負債比率の変化（財務悪化の proxy）}



\[
\mu_3 = - \Delta \left(\frac{\text{負債}}{\text{総資産}}\right)
\]



負債比率の上昇は収益性の悪化を示唆する。


\subsection{本研究で採用する $\mu$}

本研究では，データ可用性・理論的整合性・倒産企業の財務特性を踏まえ，
純資産の増減率 $\mu_1$ を主要指標として採用する。

以上より，収益性 $\mu$ は，
Merton/B\&S の理論構造における期待成長率の proxy として，
また新電力企業の基礎的収益力を表す実証変数として，
$Z_{\text{lin}}$ の構成要素に位置づけられる。


% ---------------------------------------------
% 5.3 外部ショック感応度 σ の定義と採用根拠（修正版）
% ---------------------------------------------

\section{外部ショック感応度 $\sigma$ の定義と採用根拠}

本節では，潜在倒産距離 $Z_{\text{lin}}$ の構成要素のうち，
外部ショック感応度 $\sigma$ を実証的にどのように測定するかを定義し，
その採用根拠を Merton/B\&S の理論構造および新電力市場の制度的特徴に基づき整理する。

\subsection{理論的対応：Merton/B\&S における $\sigma$ の位置づけ}

Merton モデルにおける $\sigma_V$ は企業価値のボラティリティであり，
外部ショックが企業価値に与える影響の大きさを表す。
すなわち，$\sigma_V$ が大きい企業ほど外部環境の変動が
企業価値に強く反映され，倒産確率が高まる。

Bharath \& Shumway（2008）（以下 B\&S）は，
$\sigma_V$ を株価ボラティリティによって proxy 化し，



\[
\sigma_V \approx \hat{\sigma}_{\text{market}}
\]



とすることで，外部ショックへの脆弱性を実証的に扱える形に変換した。

本研究もこの proxy 化の思想を継承しつつ，
新電力企業の制度的特徴に適合する形で $\sigma$ を定義する。

\subsection{新電力市場における外部ショック感応度の意味}

新電力企業の外部ショック感応度は，
JEPX 価格急騰や需給調整金（インバランス）の変動といった外部環境の変化が
企業の損益にどの程度影響するかを表す。

具体的には，

\begin{itemize}
    \item スポット市場への依存度が高い企業ほど，JEPX 価格急騰の影響を強く受ける
    \item 調達構造が固定費化されていない企業ほど，外部ショックが損益に直結する
    \item 外部ショック感応度は倒産距離を押し下げる（$b>0$）
\end{itemize}

という特徴を持つ。

\subsection{本来の定義：企業価値ボラティリティ（理論上の $\sigma$）}

理論的には，外部ショック感応度 $\sigma$ は，
企業価値 $V_t$ のボラティリティ $\sigma_V$ によって定義される。



\[
\sigma_{\text{theory}} = \sigma_V
\]



しかし，新電力企業の多くは非上場であり，
株価データが存在しないため，
B\&S が用いた市場ベースのボラティリティ proxy を利用できない。

したがって，$\sigma_{\text{theory}}$ を
実証的に測定することは困難である。

\subsection{実証上の定義：外部ショック感応度の proxy}

この制約を踏まえ，本研究では，
外部ショック感応度 $\sigma$ を以下の指標によって proxy 化する。

\noindent\textbf{（1）スポット依存度（最も理論的整合性が高い）}



\[
\sigma_1 = \frac{\text{スポット調達量}}{\text{総調達量}}
\]



スポット依存度は，JEPX 価格急騰が損益に与える影響を最も直接的に反映する。

\noindent\textbf{（2）JEPX 価格ボラティリティ（市場ショックの proxy）}



\[
\sigma_2 = \text{Std}(\text{JEPXスポット価格})
\]



市場全体のショック強度を表し，
外部環境の変動を企業横断的に捉える。

\noindent\textbf{（3）インバランス単価の変動係数（制度ショックの proxy）}



\[
\sigma_3 = \text{CV}(\text{インバランス単価})
\]



制度的ショックの強さを反映し，
倒産の直接原因となる負担増を捉える。

\subsection{本研究で採用する $\sigma$}

本研究では，新電力企業の倒産構造が
「スポット依存度 × 市場ショック」の掛け算的構造を持つことを踏まえ，
スポット依存度 $\sigma_1$ を主要指標として採用し，
補完的に JEPX 価格ボラティリティ $\sigma_2$ を併用する。

以上より，外部ショック感応度 $\sigma$ は，
Merton/B\&S の理論構造におけるボラティリティの proxy として，
また新電力企業の外部ショック脆弱性を表す実証変数として，
$Z_{\text{lin}}$ の構成要素に位置づけられる。


% ---------------------------------------------
% 5.4 支援能力 S の定義と採用根拠（修正版）
% ---------------------------------------------

\section{支援能力 $S$ の定義と採用根拠}

本節では，潜在倒産距離 $Z_{\text{lin}}$ の構成要素のうち，
支援能力 $S$ を実証的にどのように測定するかを定義し，
その採用根拠を Merton/B\&S の理論構造および新電力市場の制度的特徴に基づき整理する。

\subsection{理論的対応：Merton/B\&S には存在しない新規要素}

支援能力 $S$ は，Merton モデルにも B\&S にも明示的には存在しない，
本研究独自の拡張要素である。

Merton/B\&S の枠組みでは，倒産確率は



\[
\mu,\ \sigma_V,\ \ln(V/F)
\]



といった企業内部の財務構造に基づいて決まる。
しかし，新電力企業の倒産構造は，
\textbf{親会社・自治体・金融機関による外生的な支援の有無が倒産回避に決定的な影響を与える}
という点で，従来の倒産モデルとは本質的に異なる。

したがって，本研究では Merton/B\&S の構造を拡張し，
支援能力 $S$ を倒産距離の構成要素として導入する。

\subsection{新電力市場における支援能力の意味}

新電力企業は，外部ショックが発生した際に，
以下のような外部支援によって倒産を回避するケースが多い。

\begin{itemize}
    \item 親会社による資金注入・信用補完
    \item 自治体による財政支援・補助金・緊急融資
    \item 金融機関による追加融資・返済猶予
\end{itemize}

これらの支援は企業内部の財務指標では捉えられない
\textbf{外生的な倒産回避メカニズム}であり，
倒産距離を押し上げる効果を持つ（$c>0$）。

特に新電力市場では，

\begin{itemize}
    \item 親会社が大企業である場合の救済可能性
    \item 自治体新電力に対する政策的保護
    \item 金融機関の継続支援
\end{itemize}

が倒産リスクを大きく左右するため，
支援能力 $S$ をモデルに含めることは理論的にも実務的にも不可欠である。

\subsection{データ制約：支援の強度を連続変数として測定できない}

支援能力は本質的に「外生的・制度的」な要因であり，
その強度を定量化する客観的データは存在しない。

\begin{itemize}
    \item 親会社支援の金額は非公開であることが多い
    \item 自治体支援は制度ごとに形式が異なり，定量化が困難
    \item 金融機関の融資姿勢は公開情報からは把握できない
\end{itemize}

したがって，支援能力を連続変数として測定することは現実的ではなく，
\textbf{支援の有無を二値化（dummy）する}ことが最も妥当である。

\subsection{実証上の定義：支援能力の二値化（dummy）}

本研究では，支援能力 $S$ を以下の二値変数として定義する。



\[
S =
\begin{cases}
1 & \text{親会社・自治体・金融機関のいずれかから支援が確認できる場合} \\
0 & \text{支援が確認できない場合}
\end{cases}
\]



二値化する理由は次の通りである。

\begin{itemize}
    \item 支援の有無が倒産回避に与える影響は「閾値的」であり，連続変数化が難しい
    \item 支援の強度を定量化する客観的データが存在しない
    \item 倒産研究では支援の有無を dummy で扱うことが一般的である
\end{itemize}

補足的に，支援の種類を区別することも可能である。

\begin{itemize}
    \item 親会社支援 dummy
    \item 自治体支援 dummy
    \item 金融機関支援 dummy
\end{itemize}

ただし，本研究ではサンプルサイズと識別性を考慮し，
総合的な支援能力 dummy を主要指標として採用する。

\subsection{データソース}

支援能力の測定には，以下の情報源を用いる。

\begin{itemize}
    \item 親会社の有価証券報告書（子会社支援の記載）
    \item 自治体議会資料・補助金交付資料
    \item 金融機関の支援報道・融資情報
    \item 官報（破産・民事再生の公告）
    \item 新電力企業のプレスリリース
\end{itemize}

これらの情報を総合的に判断し，
支援の有無を客観的に判定する。

以上より，支援能力 $S$ は，
Merton/B\&S の理論構造を新電力市場に適合させるための
\textbf{本研究独自の拡張要素}として，
$Z_{\text{lin}}$ の構成要素に位置づけられる。


% ---------------------------------------------
% 5.5 倒産閾値 D の定義と採用根拠（修正版）
% ---------------------------------------------

\section{倒産閾値 $D$ の定義と採用根拠}

本節では，潜在倒産距離 $Z_{\text{lin}}$ の構成要素のうち，
倒産閾値 $D$ を実証的にどのように測定するかを定義し，
その採用根拠を Merton/B\&S の理論構造および新電力市場の制度的特徴に基づき整理する。

\subsection{理論的対応：Merton の負債額 $F$ の再解釈}

Merton モデルにおける倒産閾値は負債額 $F$ であり，
企業価値 $V_T$ が $F$ を下回ると倒産が発生する。
Bharath \& Shumway（2008）（以下 B\&S）もこの構造を維持し，
帳簿価値の負債額を proxy として用いた。

すなわち，



\[
\text{Merton: } \ln(V/F), \qquad
\text{B\&S: } \ln\left(\frac{E+F}{F}\right)
\]



という形で，倒産閾値は「企業が最低限クリアすべき支払義務」として扱われる。

しかし，新電力企業において倒産を引き起こす主要因は，
固定的な負債額ではなく，



\[
\textbf{外部ショック発生時に突発的に発生する支払義務}
\]



である。

したがって，本研究では Merton/B\&S の $F$ を
新電力企業の制度的特徴に合わせて再解釈し，
倒産閾値 $D$ として定義する。

\subsection{新電力市場における倒産閾値の意味}

新電力企業の倒産は，以下のような外部ショック時の支払義務が
短期間に急増することで発生する。

\begin{itemize}
    \item インバランス料金（需給調整金）の急増
    \item 調達費用（スポット価格）の急騰
    \item 制度コスト（容量市場・非化石価値取引等）の上昇
\end{itemize}

これらは企業の財務構造とは独立に発生し，
短期間で資金繰りを圧迫する。

したがって，倒産閾値 $D$ は



\[
\textbf{外部ショック時に企業が最低限支払わなければならないキャッシュアウト}
\]



として定義され，$Z_{\text{lin}}$ において倒産距離を押し下げる要因（$d>0$）となる。

\subsection{データ制約：外部ショック時の支払義務は直接観測できない}

外部ショック時の支払義務は，以下の理由により直接観測できない。

\begin{itemize}
    \item インバランス量は企業別に公開されていない
    \item 調達構造（長期契約比率・FIT 比率）は非公開であることが多い
    \item 制度コストの負担額は企業ごとに異なり，公開情報からは推計が必要
\end{itemize}

したがって，倒産閾値 $D$ を直接測定することは困難であり，
\textbf{外部ショック時の支払義務を proxy によって定量化する必要がある}。

\subsection{実証上の定義：倒産閾値 $D$ の proxy}

本研究では，倒産閾値 $D$ を以下の 3 つの要素の合算として proxy 化する。

\noindent\textbf{（1）最大インバランス負担額（最も倒産に直結）}



\[
D_{\text{imbalance}}
= \max_{t \in \text{shock}} 
\left( \text{インバランス単価}_t \times \text{インバランス量}_t \right)
\]



\noindent\textbf{（2）調達費用急騰幅（市場ショックの proxy）}



\[
D_{\text{procurement}}
= \max_{t \in \text{shock}}
\left( \text{スポット価格}_t - \text{平常時価格} \right)
\]



\noindent\textbf{（3）制度コスト上昇額（制度ショックの proxy）}



\[
D_{\text{institution}}
= \text{容量市場負担} + \text{非化石価値取引負担}
\]



\subsection{本研究で採用する $D$}

本研究では，新電力企業の倒産が
「インバランス負担 × 市場ショック × 制度コスト」
の掛け算的構造を持つことを踏まえ，



\[
D = D_{\text{imbalance}} + D_{\text{procurement}} + D_{\text{institution}}
\]



を倒産閾値として採用する。

以上より，倒産閾値 $D$ は，
Merton/B\&S の負債額 $F$ を新電力市場に適合させた
\textbf{動的な倒産閾値}として，
$Z_{\text{lin}}$ の構成要素に位置づけられる。
